\section*{Chapitre \ref{ch:qm:robust-statistics-and-heavy-tails} --- Statistiques robustes et queues épaisses}

\begin{solution}{exo:qm:robust-statistics-and-heavy-tails:1}
10~\% pour la \omterm{def:qm:robust-statistics-and-heavy-tails:winsor}{moyenne tronquée} (elle bascule si plus que la part tronquée se trouve d'un seul côté)~; 25~\% pour l'écart interquartile, qui ne dépend que
des quartiles.
\end{solution}

\begin{solution}{exo:qm:robust-statistics-and-heavy-tails:2}
Les moments d'ordre inférieur à 3 existent~: la moyenne et la variance. Le kurtosis, moment d'ordre 4, n'existe pas~; le moment d'ordre 3 est le cas limite
et il est infini pour une queue de Pareto exacte, car $\int^\infty x^2 \cdot x^{-4}\,x\,dx$ diverge logarithmiquement.
\end{solution}

\begin{solution}{exo:qm:robust-statistics-and-heavy-tails:3}
$\rho_S = \frac6\pi\arcsin0.21 = 0.404$~; $\tau = \frac2\pi\arcsin0.42 = 0.276$.
\end{solution}

\begin{solution}{exo:qm:robust-statistics-and-heavy-tails:4}
$\E[\psi_c'(Z)] = 2\Phi(c) - 1 = 0.8214$. $\E[\psi_c(Z)^2] = \E[Z^2; |Z| \le c] + c^2\P(|Z| > c) = (2\Phi(c) - 1) - 2c\varphi(c) + 2c^2(1 - \Phi(c)) =
0.7102$. La variance asymptotique est $\E[\psi^2]/\E[\psi']^2$, donc l'efficacité relative à la moyenne vaut $0.8214^2/0.7102 = 0.950$.
\end{solution}

\begin{solution}{exo:qm:robust-statistics-and-heavy-tails:5}
Au-dessus de $x_0 = X_{(k+1)}$, la densité de Pareto vaut $\alpha x_0^\alpha x^{-\alpha-1}$. La log-vraisemblance des $k$ plus grandes valeurs est $k\ln\alpha + k\alpha\ln x_0 - (\alpha +
1)\sum_{i \le k}\ln X_{(i)}$~; sa dérivée $k/\alpha - \sum_i(\ln X_{(i)} - \ln x_0)$ s'annule en $\hat\alpha = k/\sum_i\ln(X_{(i)}/x_0)$, soit
l'\omterm{def:qm:robust-statistics-and-heavy-tails:hill}{estimateur de Hill}.
\end{solution}

\begin{solution}{exo:qm:robust-statistics-and-heavy-tails:6}
$0.93 + \frac{0.34}{0.069}\bigl((0.05/0.0001)^{0.069} - 1\bigr) = 3.59\%$ (avec l'ajustement non arrondi~; $N_u/n = 355/7\,098 = 0.05$). Deux jours en 27 ans
ont perdu davantage (3.7~\% et 4.7~\%), à peu près ce que l'ajustement implique (0.7 attendu)~; le peu d'observations explique
pourquoi l'incertitude de l'extrapolation est grande.
\end{solution}

\begin{solution}{exo:qm:robust-statistics-and-heavy-tails:7}
Au-dessus des quantiles à 90~\%, 95~\%, 97.5~\% et 99~\% (710, 355, 178 et 71 excès)~: $\xi = 0.033$, $0.069$, $0.098$, $0.212$~; la perte
d'un jour sur mille~: 2.45~\%, 2.47~\%, 2.49~\%, 2.43~\%. La forme fluctue, le quantile dans les données, lui, ne fluctue pas.
\end{solution}

\begin{solution}{exo:qm:robust-statistics-and-heavy-tails:8}
Sur 250 jours, la \omterm{def:qm:robust-statistics-and-heavy-tails:winsor}{winsorisation} à 1~\% plafonne les deux ou trois plus fortes variations de chaque côté. Elle neutralise par hasard le seul mauvais relevé du chapitre
(l'écart type corrompu devient 0.341~\% contre 0.339~\% propre), mais elle plafonne aussi les grands mouvements authentiques chaque jour, ce qui biaise la
volatilité propre vers le bas (0.328~\% contre 0.339~\%), surtout dans les périodes de tension qui comptent~; et trois mauvais relevés en un an passent à travers. Par ailleurs,
l'écart type des données winsorisées exige aussi une correction de cohérence. Filtrer les relevés en amont, et utiliser une échelle robuste dès la conception.
\end{solution}

\begin{solution}{pb:qm:robust-statistics-and-heavy-tails:1}
\textbf{1.} $\ln(1.5172/1.1572) = 0.271$~: les rendements des 24 et 25 mars 2026 deviennent environ $+27\%$ et $-27\%$.
\textbf{2.} Propres~: 0.339~\%, 0.275~\%, 0.300~\%. Corrompues~: 2.434~\%, 0.279~\%, 0.305~\%.
\textbf{3.} La moyenne ($-0.011\%$) et la médiane ($-0.030\%$) ne changent pas, la première par chance, car les deux mauvais rendements se compensent~; la
position de Huber varie à peine ($-0.01953\%$ à $-0.01948\%$), son influence bornée jouant pleinement.
\textbf{4.} Dix-sept mouvements au-delà de quatre écarts types (six pertes, onze gains), contre 0.45 attendu sous une loi normale.
\textbf{5.} 6.3~; un moment d'ordre 4 d'une loi dont l'\omterm{def:qm:robust-statistics-and-heavy-tails:heavy}{indice de queue} est proche de 4 est à peine fini, et une poignée de jours décident de l'estimation.
\textbf{6.} Environ 3.9 sur un palier pour $k$ entre 100 et 200 plus grandes pertes (4.9 pour $k = 50$, glissant jusqu'à 2.6 pour $k = 700$).
\textbf{7.} $u = 0.93\%$, 355 excès, $\xi = 0.069$, $\beta = 0.34\%$.
\textbf{8.} $\nu = 4.7$, échelle 0.44~\%.
\textbf{9.} Normale 1.80~\%, Student $t$ 2.74~\%, GPD 2.47~\%~; empirique 2.27~\%.
\textbf{10.} 35 jours au-delà du quantile normal, 6 au-delà du quantile GPD.
\textbf{11.} Normale 2.27~\% (+27~\%), Student $t$ 2.91~\% (+6~\%), GPD 2.85~\% (+15~\%).
\textbf{12.} Le quantile normal vaut l'écart type multiplié par 3.09, et l'écart type a une \omterm{def:qm:robust-statistics-and-heavy-tails:influence}{fonction d'influence} non bornée~; une
paire de rendements de $\pm 27\%$ la fait passer de 0.58~\% à 0.74~\%.
\textbf{13.} Il passe de 0.07 à 0.22~: la mauvaise perte est de loin le plus grand excès, et le paramètre de forme est ajusté surtout par les plus grands
excès.
\textbf{14.} Sa vraisemblance pondère une observation par $(\nu + 1)/(\nu + z^2)$, poids qui décroît avec la taille du résidu~: la vraisemblance de la loi de Student-$t$ est elle-même un
\omterm{def:qm:estimation:m}{M-estimateur} robuste.
\textbf{15.} En partie~: un quantile normal fondé sur l'\omterm{def:qm:robust-statistics-and-heavy-tails:huber}{écart absolu médian} ignore le mauvais relevé, mais il est encore plus éloigné de la vraie queue (la forme normale est fausse)~; une
échelle robuste ne rend pas un modèle à queue fine apte à une variable à queue épaisse.
\textbf{16.} Le quantile GPD ou de Student-$t$, environ 2.5~\% à 2.7~\%, présenté avec son seuil et sa sensibilité, jamais le quantile normal de 1.8~\%.
\textbf{17.} À l'ingestion~: un filtre qui signale un relevé dont le rendement dépasse de nombreuses échelles robustes et s'inverse le lendemain, avant que la
moindre statistique soit calculée.
\textbf{18.} Les pertes conjointes sont plus fréquentes qu'une \omterm{def:qm:robust-statistics-and-heavy-tails:copula}{copule} gaussienne ne le suppose (32~\% contre 20~\% dans la queue conjointe à 5~\%), si bien qu'une limite sur la somme des
deux expositions, calculée avec une corrélation de 0.42, sous-estime leur perte conjointe en situation de stress.
\textbf{19.} Résultat baptisé~: \emph{la mauvaise cotation}~: la perte EUR/USD d'un jour sur mille vaut 1.80~\% (normale), 2.74~\% (Student $t$) et 2.47~\%
(GPD au-delà du quantile à 95~\%)~; un seul relevé à chiffres permutés les fait monter de 27~\%, 6~\% et 15~\%.
\textbf{20.} Les estimations du centre et de l'échelle du corps de la distribution peuvent être rendues robustes~; celles de la queue ne le peuvent pas, car dans la queue chaque
observation est une information, si bien que la protection doit venir du nettoyage des données.
\end{solution}

\begin{solution}{iq:qm:robust-statistics-and-heavy-tails:1}
Les données~: un mauvais relevé ou un cours obsolète puis corrigé produit une paire de grands rendements opposés. Regarder les plus grands rendements de la fenêtre,
comparer avec une autre source, et comparer l'écart type à l'\omterm{def:qm:robust-statistics-and-heavy-tails:huber}{écart absolu médian}~: un écart septuple entre les deux relève d'un problème de données, et non du marché.
\iqlookfor{Les données avant le modèle~; la signature des retournements appariés~; une échelle robuste comme diagnostic.}
\end{solution}

\begin{solution}{iq:qm:robust-statistics-and-heavy-tails:2}
La plus petite fraction de contamination arbitraire qui peut entraîner l'\omterm{def:qm:estimation:estimator}{estimateur} arbitrairement loin. Moyenne~: $1/n$, qui tend vers zéro. Médiane~: un demi.
\omterm{def:qm:robust-statistics-and-heavy-tails:huber}{Écart absolu médian}~: un demi.
\iqlookfor{La définition, et le contraste entre zéro et un demi.}
\end{solution}

\begin{solution}{iq:qm:robust-statistics-and-heavy-tails:3}
Les moments d'ordre inférieur à 3 existent, ceux d'ordre supérieur à 3 n'existent pas~: variance finie, kurtosis infini. Le kurtosis empirique ne converge alors pas~; il croît avec
la taille de l'échantillon, et il est dominé par la plus grande observation~: il ne doit donc pas servir de paramètre.
\iqlookfor{La règle des moments, et la conséquence pour les moments empiriques.}
\end{solution}

\begin{solution}{iq:qm:robust-statistics-and-heavy-tails:4}
Environ six observations se situent au-delà du quantile à 99.9~\% en 25 ans~: le quantile empirique est bruité, un ajustement normal est fortement biaisé. Ajuster une
\omterm{def:qm:robust-statistics-and-heavy-tails:gpd}{loi de Pareto généralisée} au-delà d'un seuil élevé (\omterm{def:qm:robust-statistics-and-heavy-tails:gpd}{méthode des excès au-delà d'un seuil}, POT), vérifier la stabilité selon les seuils, comparer à un ajustement de Student-$t$,
traiter le regroupement de volatilité en ajustant la loi aux rendements normalisés par une prévision de volatilité, et rapporter l'incertitude.
\iqlookfor{\omterm{def:qm:robust-statistics-and-heavy-tails:gev}{Théorie des valeurs extrêmes} avec diagnostics de seuil, et prise en compte de la dépendance.}
\end{solution}

\begin{solution}{iq:qm:robust-statistics-and-heavy-tails:5}
Les \omterm{def:qm:robust-statistics-and-heavy-tails:rank}{corrélations de rang} sont invariantes par transformations croissantes et bornées en influence, si bien que les valeurs aberrantes et les queues épaisses ne les dominent pas. Ni
la corrélation de Pearson ni la \omterm{def:qm:robust-statistics-and-heavy-tails:rank}{corrélation de rang} ne décrivent les queues~: deux couples de même \omterm{def:qm:robust-statistics-and-heavy-tails:rank}{corrélation de rang} peuvent avoir des probabilités de krach conjoint
très différentes.
\iqlookfor{Invariance et robustesse des rangs, et dépendance de queue comme quantité distincte.}
\end{solution}

\begin{solution}{iq:qm:robust-statistics-and-heavy-tails:6}
Toute loi conjointe s'écrit à partir de ses marges et d'une \omterm{def:qm:robust-statistics-and-heavy-tails:copula}{copule}, unique pour des marges continues, si bien que la dépendance peut être modélisée séparément des marges. La
\omterm{def:qm:robust-statistics-and-heavy-tails:copula}{copule} gaussienne a une dépendance de queue nulle pour toute corrélation inférieure à un~: les mouvements conjoints extrêmes deviennent asymptotiquement indépendants, alors que sur les marchés
ils se regroupent, ce que permet une \omterm{def:qm:robust-statistics-and-heavy-tails:copula}{copule} de Student-$t$.
\iqlookfor{La décomposition, et la dépendance de queue nulle ou positive.}
\end{solution}
